% Part: computability
% Chapter: computability-theory
% Section: non-comp-set

\documentclass[../../../include/open-logic-section]{subfiles}

\begin{document}

\olfileid{cmp}{thy}{ncp}
\olsection{संगणनक्षम नसलेले संच अस्तित्वात आहेत}


प्रत्येक संगणनक्षम संच संगणनक्षमपणे प्रगणनीय असतो, हे आपण वर पाहिले.
याचे उलटही खरे आहे का? पुढील प्रमेय दाखवते की सर्वसाधारणपणे तसे नाही.

\begin{thm}\ollabel{thm:K-0}
$K_0$ हा $\Setabs{\tuple{e, x}}{\cfind{e}(x) \fdefined}$ हा
संच असू द्या. मग $K_0$ संगणनक्षमपणे प्रगणनीय आहे, पण संगणनक्षम नाही.
\end{thm}

\begin{proof}
$K_0$ संगणनक्षमपणे प्रगणनीय आहे हे पाहण्यासाठी, तो पुढीलप्रमाणे
परिभाषित फलन~$f$ चे परिभाषाक्षेत्र आहे हे लक्षात घ्या:
\[
f(z) = \umin{y}{(\len{z} = 2 \land T((z)_0, (z)_1, y))}.
\]
कारण $\cfind{e}(x)$ परिभाषित असेल, तर $f(\tuple{e, x})$
थांबणारी संगणनक्रमिका शोधते; $\cfind{e}(x)$ अपरिभाषित असेल,
तर $f(\tuple{e, x})$ देखील अपरिभाषित असते; आणि $z$ हा
जोडीचा संकेतांकसुद्धा नसेल, तर $f(z)$ देखील अपरिभाषित असते.

$K_0$ संगणनक्षम नाही, ही बाब म्हणजे थांबण्याच्या समस्येचीच
अनिर्णेयता, \olref[hlt]{thm:halting-problem}.
\end{proof}

संच $K_0$ म्हणजे $\cfind{e}(x) \fdefined$ असा गुणधर्म असलेल्या
$\tuple{e,x}$ जोड्यांचा संच. म्हणजे $\tuple{e,x} \in K_0$
असेल तेव्हाच $\cfind{e}$ हे आदान~$x$ वर परिभाषित असते
(संगणन थांबते). म्हणून त्याला ``थांबणारा संच'' असेही म्हणतात.
$K = \Setabs{e}{\cfind{e}(e) \fdefined}$ हा संच
``स्वतःच्या निर्देशांकावर थांबणारा संच'' आहे. तो अनिर्णेय संचाचे
मानक उदाहरण म्हणून वारंवार वापरला जातो.

\begin{thm}\ollabel{thm:K}
स्वतःच्या निर्देशांकावर थांबणारा संच $K = \Setabs{e}{\cfind{e}(e) \fdefined}$
हा !!{c.e.} आहे, पण निर्णेय नाही.
\end{thm}

\begin{proof}
  $K$ निर्णेय आहे, म्हणजे त्याचे जातिबोधक फलन
  $\Char{K}$ संगणनक्षम आहे, असे समजा. पुढीलप्रमाणे घ्या:
  \[d(e) = \begin{cases}
  1 & \text{जर\/ $\Char{K}(e) = 0$ असेल तर}\\
  \fundefined & \text{अन्यथा.}
  \end{cases}
  \]
  $k$ हा~$d$ चा निर्देशांक असू द्या, म्हणजे $d \simeq \cfind{k}$.
  मग $d(k)
  \simeq \cfind{k}(k)$. पण $d(k)
  \fdefined$ असेल तेव्हाच $\cfind{k}(k) \fundefined$, हे~$d$
  च्या व्याख्येवरून मिळते; त्यामुळे विरोध होतो.

  $K$ हे $f(x) = \umin{y}{T(x,x,y)}$ चे परिभाषाक्षेत्र आहे,
  म्हणून ते !!{c.e.} आहे.
\end{proof}

\end{document}
